% \begin{document}

\chapter{Lambda Calculus}

Lambda calculus (or $\lambda$-calculus) was introduced by Alonzo Church in the 1930s, and developed
with his students Stephen Kleene and J.\ Barkley Rosser, to describe functions in an unambiguous
and compact manner. Many real languages are based on the lambda calculus, including Lisp, Scheme,
Haskell, and ML. A key characteristic of these languages is that functions are values, just like
integers and booleans are values: functions can be used as arguments to functions, and can be
returned from functions.

The name ``lambda calculus'' comes from the use of the Greek letter lambda ($\lambda$) in function
definitions. (The letter lambda has no significance.) ``Calculus'' means a method of calculating
by the symbolic manipulation of expressions.

Intuitively, a function is a rule for determining a value from an argument. Some examples of
functions in mathematics are
\begin{align*}
  f(x)&=x^3\\
  g(y)&=y^3-2y^2+5y-6.
\end{align*}

\section{Syntax}

The pure $\lambda$-calculus contains just function definitions (called \emph{abstractions}),
variables, and function application (i.e., applying a function to an argument). If we add
additional data types and operations (such as integers and addition), we have an \emph{applied}
$\lambda$-calculus. In the following text, we will sometimes assume that we have integers and
addition in order to give more intuitive examples.

The syntax of the pure $\lambda$-calculus is defined as follows.
\[
  \begin{array}{lcp{0.3\textwidth}r}
    e\ ::=\ x & & & \text{variable}\\
    \quad\mid\ \lambda x.\,e & & & \text{abstraction}\\
    \quad\mid\ e_1\ e_2 & & & \text{application}
  \end{array}
\]
An abstraction $\lambda x.\,e$ is a function: variable $x$ is the \emph{argument}, and expression
$e$ is the \emph{body} of the function. Note that the function $\lambda x.\,e$ doesn't have a name.
Assuming we have integers and arithmetic operations, the expression $\lambda x.\,x^2$ is a function
that takes an argument $x$ and returns the square of $x$.

An application $e_1\ e_2$ requires that $e_1$ is (or evaluates to) a function, and then applies the
function to the expression $e_2$. For example, $(\lambda x.\,x^2)\ 5$ is, intuitively, equal to
$25$, the result of applying the squaring function $\lambda x.\,x^2$ to $5$.

Here are some examples of lambda calculus expressions.\footnote{Lambda expressions extend as far to
the right as possible. For example $\lambda x.\,x\ \lambda y.\,y$ is the same as
$\lambda x.\,x\ (\lambda y.\,y)$, and is not the same as $(\lambda x.\,x)\ (\lambda y.\,y)$.
Application is left associative. For example $e_1\ e_2\ e_3$ is the same as $(e_1\ e_2)\ e_3$. In
general, use parentheses to make the parsing of a lambda expression clear if you are in doubt.}

\begin{center}
\begin{tabular}{ll}
  $\lambda x.\,x$ & a lambda abstraction called the \emph{identity function}\\
  $\lambda x.\,f\ (g\ x)$ & another abstraction\\
  $(\lambda x.\,x)\ 42$ & an application\\
  $\lambda y.\,\lambda x.\,x$ & an abstraction that ignores its argument and returns the identity function
\end{tabular}
\end{center}

\subsection{Variable binding and $\alpha$-equivalence}

An occurrence of a variable in an expression is either bound or free. An occurrence of a variable
$x$ in a term is \emph{bound} if there is an enclosing $\lambda x.\,e$; otherwise, it is
\emph{free}. A \emph{closed} term is one in which all identifiers are bound.

Consider the following term:
\[
  \lambda x.\,(x\ (\lambda y.\,y\ a)\ x)\ y
\]
Both occurrences of $x$ are bound, the first occurrence of $y$ is bound, the $a$ is free, and the
last $y$ is also free, since it is outside the scope of the $\lambda y$.

If a program has some variables that are free, then you do not have a complete program as you do not
know what to do with the free variables. Hence, a well formed program in lambda calculus is a
closed term.

The symbol $\lambda$ is a \emph{binding operator}, as it binds a variable within some scope (i.e.,
some part of the expression): variable $x$ is bound in $e$ in the expression $\lambda x.\,e$.

The name of bound variables is not important. Consider the mathematical integrals
$\int_0^7 x^2\,dx$ and $\int_0^7 y^2\,dy$. They describe the same integral, even though one uses
variable $x$ and the other uses variable $y$ in their definition. The meaning of these integrals is
the same: the bound variable is just a placeholder. In the same way, we can change the name of
bound variables without changing the meaning of functions. Thus $\lambda x.\,x$ is the same
function as $\lambda y.\,y$. Expressions $e_1$ and $e_2$ that differ only in the name of bound
variables are called \emph{$\alpha$-equivalent}, sometimes written $e_1\equiv_\alpha e_2$.

\subsection{Higher-order functions}

In lambda calculus, functions are values: functions can take functions as arguments and return
functions as results. In the pure lambda calculus, every value is a function, and every closed
program that terminates produces a function!

For example, the following function takes a function $f$ as an argument, and applies it to the
value $42$.
\[
  \lambda f.\,f\ 42
\]
This function takes an argument $v$ and returns a function that applies its own argument (a
function) to $v$.
\[
  \lambda v.\,\lambda f.\,(f\ v)
\]

\section{Semantics}

\subsection{$\beta$-equivalence}

Application $(\lambda x.\,e_1)\ e_2$ applies the function $\lambda x.\,e_1$ to $e_2$. In some ways,
we would like to regard the expression $(\lambda x.\,e_1)\ e_2$ as equivalent to the expression
$e_1$ where every (free) occurrence of $x$ is replaced with $e_2$. For example, we would like to
regard $(\lambda x.\,x^2)\ 5$ as equivalent to $5^2$.\footnote{We write $\sub{e_1}{e_2}{x}$ to mean
expression $e_1$ with all free occurrences of $x$ replaced with $e_2$. There are several different
notations to express this substitution, including $[x\mapsto e_2]e_1$ (used by Pierce),
$[e_2/x]e_1$ (used by Mitchell), and $e_1[e_2/x]$ (used by Winskel).}

Using our notation, we would like expressions $(\lambda x.\,e_1)\ e_2$ and $\sub{e_1}{e_2}{x}$ to be
equivalent. This equivalence, between $(\lambda x.\,e_1)\ e_2$ and $\sub{e_1}{e_2}{x}$, is called
\emph{$\beta$-equivalence}. Rewriting $(\lambda x.\,e_1)\ e_2$ into $\sub{e_1}{e_2}{x}$ is called a
\emph{$\beta$-reduction}. Given a lambda calculus expression, we may, in general, be able to
perform $\beta$-reductions. This corresponds to executing a lambda calculus expression.

There may be more than one possible way to $\beta$-reduce an expression. Consider, for example,
$(\lambda x.\,x+x)\ ((\lambda y.\,y)\ 5)$. We could use $\beta$-reduction to get either
$((\lambda y.\,y)\ 5)+((\lambda y.\,y)\ 5)$ or $(\lambda x.\,x+x)\ 5$. The order in which we
perform $\beta$-reductions results in different semantics for the lambda calculus.

\subsection{Call-by-value}

Call-by-value (or CBV) semantics makes sure that functions are only called on values. That is,
given an application $(\lambda x.\,e_1)\ e_2$, CBV semantics makes sure that $e_2$ is a value
before calling the function.

So, what is a value? In the pure lambda calculus, any abstraction is a value. Remember, an
abstraction $\lambda x.\,e$ is a function; in the pure lambda calculus, the only values are
functions. In an applied lambda calculus with integers and arithmetic operations, values also
include integers. Intuitively, a value is an expression that can not be
reduced/executed/simplified any further.

We can give small step operational semantics for call-by-value execution of the lambda calculus.
Here, $v$ can be instantiated with any value (e.g., a function).
\[
  \rul{e_1\to e_1'}{e_1\ e_2\to e_1'\ e_2}
  \qquad
  \rul{e\to e'}{v\ e\to v\ e'}
  \qquad
  \infr{}{(\lambda x.\,e)\ v\to\sub{e}{v}{x}}{$\beta$-Reduction}
\]
We can see from these rules that, given an application $e_1\ e_2$, we first evaluate $e_1$ until it
is a value, then we evaluate $e_2$ until it is a value, and then we apply the function to the
value---a $\beta$-reduction.

Let's consider some examples. (These examples use an applied lambda calculus that also includes
reduction rules for arithmetic expressions.)
\begin{align*}
  (\lambda x.\,\lambda y.\,y\ x)\ (5+2)\ \lambda x.\,x+1
    &\to(\lambda x.\,\lambda y.\,y\ x)\ 7\ \lambda x.\,x+1\\
    &\to(\lambda y.\,y\ 7)\ \lambda x.\,x+1\\
    &\to(\lambda x.\,x+1)\ 7\\
    &\to 7+1\\
    &\to 8
\end{align*}
\begin{align*}
  (\lambda f.\,f\ 7)\ ((\lambda x.\,x\ x)\ \lambda y.\,y)
    &\to(\lambda f.\,f\ 7)\ ((\lambda y.\,y)\ (\lambda y.\,y))\\
    &\to(\lambda f.\,f\ 7)\ (\lambda y.\,y)\\
    &\to(\lambda y.\,y)\ 7\\
    &\to 7
\end{align*}

\subsection{Call-by-name}

Call-by-name (or CBN) semantics applies the function as soon as possible. The small step
operational semantics are a little simpler, as they do not need to ensure that the expression to
which a function is applied is a value.
\[
  \rul{e_1\to e_1'}{e_1\ e_2\to e_1'\ e_2}
  \qquad
  \infr{}{(\lambda x.\,e_1)\ e_2\to\sub{e_1}{e_2}{x}}{$\beta$-Reduction}
\]
Let's consider the same examples we used for CBV.
\begin{align*}
  (\lambda x.\,\lambda y.\,y\ x)\ (5+2)\ \lambda x.\,x+1
    &\to(\lambda y.\,y\ (5+2))\ \lambda x.\,x+1\\
    &\to(\lambda x.\,x+1)\ (5+2)\\
    &\to(5+2)+1\\
    &\to 7+1\\
    &\to 8
\end{align*}
\begin{align*}
  (\lambda f.\,f\ 7)\ ((\lambda x.\,x\ x)\ \lambda y.\,y)
    &\to((\lambda x.\,x\ x)\ \lambda y.\,y)\ 7\\
    &\to((\lambda y.\,y)\ (\lambda y.\,y))\ 7\\
    &\to(\lambda y.\,y)\ 7\\
    &\to 7
\end{align*}
Note that the answers are the same, but the order of evaluation is different. (Later we will see
languages where the order of evaluation is important, and may result in different answers.)

\section{Lambda calculus evaluation}

There are many different evaluation strategies for the $\lambda$-calculus. The most permissive is
\emph{full $\beta$ reduction}, which allows any \emph{redex}---i.e., any expression of the form
$(\lambda x.\,e_1)\ e_2$---to step to $\sub{e_1}{e_2}{x}$ at any time. It is defined formally by
the following small-step operational semantics rules:
\[
  \rul{e_1\to e_1'}{e_1\ e_2\to e_1'\ e_2}
  \qquad
  \rul{e_2\to e_2'}{e_1\ e_2\to e_1\ e_2'}
  \qquad
  \rul{e_1\to e_1'}{\lambda x.\,e_1\to\lambda x.\,e_1'}
  \qquad
  \infr{}{(\lambda x.\,e_1)\ e_2\to\sub{e_1}{e_2}{x}}{$\beta$}
\]
The call-by-value and call-by-name strategies from the previous section are more restrictive.
Neither allows evaluation under a $\lambda$. CBV only allows an application to reduce after its
argument has been reduced to a value (i.e., a $\lambda$-abstraction), while CBN allows an
application to reduce even when its argument is not a value. Unlike full $\beta$ reduction, both
are deterministic: every expression can take at most one step.

\section{Confluence}

It is not hard to see that the full $\beta$ reduction strategy is non-deterministic. This raises an
interesting question: do the choices made during the evaluation of an expression affect the final
result? The answer turns out to be no: full $\beta$ reduction is \emph{confluent} in the following
sense:

\medskip
\noindent\textbf{Theorem (Church--Rosser).} If $e\to^* e_1$ and $e\to^* e_2$ then there exists
$e'$ such that $e_1\to^* e'$ and $e_2\to^* e'$.

\medskip
Confluence can be depicted graphically as follows:
\begin{center}
\begin{tikzpicture}[>=Stealth]
  \node (e) at (0,2) {$e$};
  \node (e1) at (-1.6,0.7) {$e_1$};
  \node (e2) at (1.6,0.7) {$e_2$};
  \node (ep) at (0,-0.6) {$e'$};
  \draw[->] (e) -- node[above left] {$*$} (e1);
  \draw[->] (e) -- node[above right] {$*$} (e2);
  \draw[->] (e1) -- node[below left] {$*$} (ep);
  \draw[->] (e2) -- node[below right] {$*$} (ep);
\end{tikzpicture}
\end{center}
This theorem is due to Church and Rosser~\cite{church-rosser}, and confluence is often also called
the \emph{Church--Rosser property}. An important consequence is that every expression has at most
one \emph{normal form}---an expression that cannot step---up to $\alpha$-equivalence. If
$e\to^* e_1$ and $e\to^* e_2$ where $e_1$ and $e_2$ are both normal forms, then confluence gives an
$e'$ that both reduce to, which must be $e_1$ and $e_2$ themselves.

Confluence does not say that every reduction sequence reaches the normal form. For example,
$(\lambda x.\,\lambda y.\,y)\ \Omega$, where $\Omega$ is the diverging expression discussed in
Section~\ref{sec:nonterm}, can step to the normal form $\lambda y.\,y$, but it can also step to
itself forever by reducing $\Omega$.

Interestingly, the confluence property does not hold for single steps. For example, consider the
expression $e=(\lambda x.\,x\ x)\ ((\lambda y.\,y)\ z)$, which has two redexes:
\begin{center}
\begin{tikzpicture}[>=Stealth,node distance=1.2cm]
  \node (top) at (0,3.6) {$(\lambda x.\,x\ x)\ ((\lambda y.\,y)\ z)$};
  \node (a) at (-3.2,2.2) {$((\lambda y.\,y)\ z)\ ((\lambda y.\,y)\ z)$};
  \node (b) at (-3.2,0.9) {$z\ ((\lambda y.\,y)\ z)$};
  \node (c) at (3.2,2.2) {$(\lambda x.\,x\ x)\ z$};
  \node (d) at (0,-0.3) {$z\ z$};
  \draw[->] (top) -- node[above left] {outer} (a);
  \draw[->] (a) -- (b);
  \draw[->] (b) -- (d);
  \draw[->] (top) -- node[above right] {inner} (c);
  \draw[->] (c) -- (d);
\end{tikzpicture}
\end{center}
Reducing the outer redex duplicates the inner one, so the left path needs two more steps to reach
$z\ z$ while the right path needs only one. Hence there is no way to complete the diamond with
single steps on both sides, and a naive proof of confluence by induction on single steps does not
work. The standard proof, due to Tait and Martin-L\"of~\cite{barendregt}, instead uses
\emph{parallel reduction}, which contracts any set of redexes in an expression at once. Parallel
reduction does have the diamond property, and its reflexive-transitive closure coincides with
$\to^*$, so confluence follows. Takahashi~\cite{takahashi} gives a particularly short version of
this proof.

\section{Substitution}

Each of the evaluation relations for $\lambda$-calculus has a $\beta$ rule defined in terms of a
substitution operation on expressions. Because the expressions involved in the substitution may
share some variable names (and because we are working up to $\alpha$-equivalence) the definition of
this operation is slightly subtle and defining it precisely turns out to be trickier than might
first appear.

As a first attempt, consider an obvious (but incorrect) definition of the substitution operator.
Here we are substituting $e$ for $x$ in some other expression:
\begin{align*}
  \sub{y}{e}{x}&=\begin{cases} e & \text{if } y=x\\ y & \text{otherwise}\end{cases}\\
  \sub{(e_1\ e_2)}{e}{x}&=(\sub{e_1}{e}{x})\ (\sub{e_2}{e}{x})\\
  \sub{(\lambda y.\,e_1)}{e}{x}&=\lambda y.\,\sub{e_1}{e}{x}\qquad\text{where } y\neq x
\end{align*}
The intuitive idea is that the last rule relies on $\alpha$-equivalence to ``rewrite'' abstractions
that use $x$ so they do not conflict. Unfortunately, this definition produces the wrong results
when we substitute an expression with free variables under a $\lambda$. For example,
\[
  \sub{(\lambda y.\,x)}{y}{x}=(\lambda y.\,y)
\]
To fix this problem, we need to revise our definition so that when we substitute under a $\lambda$
we do not accidentally bind variables in the expression we are substituting. The following
definition correctly implements \emph{capture-avoiding substitution}:
\begin{align*}
  \sub{y}{e}{x}&=\begin{cases} e & \text{if } y=x\\ y & \text{otherwise}\end{cases}\\
  \sub{(e_1\ e_2)}{e}{x}&=(\sub{e_1}{e}{x})\ (\sub{e_2}{e}{x})\\
  \sub{(\lambda y.\,e_1)}{e}{x}&=\lambda y.\,(\sub{e_1}{e}{x})\qquad\text{where } y\neq x\text{ and } y\notin\mathrm{fv}(e)
\end{align*}
Note that in the case for $\lambda$-abstractions, we require that the bound variable $y$ be
different from the variable $x$ we are substituting for and that $y$ not appear in the free
variables of $e$, the expression we are substituting. Because we work up to $\alpha$-equivalence,
we can always pick $y$ to satisfy these side conditions. For example, to calculate
$\sub{(\lambda z.\,x\ z)}{(w\ y\ z)}{x}$ we first rewrite $\lambda z.\,x\ z$ to $\lambda u.\,x\ u$
and then apply the substitution, obtaining $\lambda u.\,(w\ y\ z)\ u$ as the result.

\section{$\lambda$-calculus encodings}

The pure $\lambda$-calculus contains only functions as values. It is not exactly easy to write
large or interesting programs in the pure $\lambda$-calculus. We can however encode objects, such
as booleans, and integers.

\subsection{Booleans}

Let us start by encoding constants and operators for booleans. That is, we want to define functions
$\tr$, $\fl$, $\AND$, $\NOT$, $\IF$, and other operators that behave as expected. For example:
\begin{align*}
  \AND\ \tr\ \fl&=\fl\\
  \NOT\ \fl&=\tr\\
  \IF\ \tr\ e_1\ e_2&=e_1\\
  \IF\ \fl\ e_1\ e_2&=e_2
\end{align*}
Let's start by defining $\tr$ and $\fl$:
\begin{align*}
  \tr&\triangleq\lambda x.\,\lambda y.\,x\\
  \fl&\triangleq\lambda x.\,\lambda y.\,y
\end{align*}
Thus, both $\tr$ and $\fl$ are functions that take two arguments; $\tr$ returns the first, and
$\fl$ returns the second. We want the function $\IF$ to behave like
\[
  \lambda b.\,\lambda t.\,\lambda f.\ \mathbf{if}\ b=\tr\ \mathbf{then}\ t\ \mathbf{else}\ f.
\]
The definitions for $\tr$ and $\fl$ make this very easy.
\[
  \IF\triangleq\lambda b.\,\lambda t.\,\lambda f.\,b\ t\ f
\]
Definitions of other operators are also straightforward.
\begin{align*}
  \NOT&\triangleq\lambda b.\,b\ \fl\ \tr\\
  \AND&\triangleq\lambda b_1.\,\lambda b_2.\,b_1\ b_2\ \fl\\
  \OR&\triangleq\lambda b_1.\,\lambda b_2.\,b_1\ \tr\ b_2
\end{align*}

\subsection{Church numerals}

Church numerals encode a number $n$ as a function that takes $f$ and $x$, and applies $f$ to $x$
$n$ times.
\begin{align*}
  0&\triangleq\lambda f.\,\lambda x.\,x\\
  1&\triangleq\lambda f.\,\lambda x.\,f\ x\\
  2&\triangleq\lambda f.\,\lambda x.\,f\ (f\ x)\\
  \SUCC&\triangleq\lambda n.\,\lambda f.\,\lambda x.\,f\ (n\ f\ x)
\end{align*}
In the definition for $\SUCC$, the expression $n\ f\ x$ applies $f$ to $x$ $n$ times (assuming that
variable $n$ is the Church encoding of the natural number $n$). We then apply $f$ to the result,
meaning that we apply $f$ to $x$ $n+1$ times.

Given the definition of $\SUCC$, we can easily define addition. Intuitively, the natural number
$n_1+n_2$ is the result of applying the successor function $n_1$ times to $n_2$.
\[
  \PLUS\triangleq\lambda n_1.\,\lambda n_2.\,n_1\ \SUCC\ n_2
\]
Similarly, multiplication adds $n_2$ to $0$, $n_1$ times, and a test for zero applies the constant
function returning $\fl$ to $\tr$ $n$ times:
\begin{align*}
  \TIMES&\triangleq\lambda n_1.\,\lambda n_2.\,n_1\ (\PLUS\ n_2)\ 0\\
  \ISZERO&\triangleq\lambda n.\,n\ (\lambda z.\,\fl)\ \tr
\end{align*}
A predecessor function $\PRED$, with $\PRED\ (n+1)=n$ and $\PRED\ 0=0$, is also definable, but it
is surprisingly tricky.

\section{Nontermination}\label{sec:nonterm}

Consider the expression $(\lambda x.\,x\ x)\ (\lambda x.\,x\ x)$, which we will refer to as
$\Omega$ for brevity. Let's try evaluating $\Omega$:
\begin{align*}
  \Omega&=(\lambda x.\,x\ x)\ (\lambda x.\,x\ x)\\
  &\to(\lambda x.\,x\ x)\ (\lambda x.\,x\ x)\\
  &=\Omega
\end{align*}
Evaluating $\Omega$ never reaches a value! It is an infinite loop!

What happens if we use $\Omega$ as an actual argument to a function? Consider this program:
\[
  (\lambda x.(\lambda y.y))\ \Omega
\]
If we use CBV semantics to evaluate the program, we must reduce $\Omega$ to a value before we can
apply the function. But $\Omega$ never evaluates to a value, so we can never apply the function.
Thus, under CBV semantics, this program does not terminate. If we use CBN semantics, we can apply
the function immediately, without needing to reduce the actual argument to a value:
\[
  (\lambda x.(\lambda y.y))\ \Omega\ \to_{\mathrm{CBN}}\ \lambda y.y
\]
This difference highlights the fact that evaluation orders, like CBV and CBN, can differ in terms of
termination: expressions that diverge under one evaluation order may terminate in another. This does
not contradict confluence: the CBV evaluation of this program simply never finishes, so it never
reaches a different result.

\section{Recursion}

$\Omega$ demonstrates that we can write nonterminating $\lambda$-terms, but so far they are not very
useful. Can we write recursive functions in the $\lambda$-calculus that do actually terminate and
produce a value, as we can in most ``real'' programming languages? It seems difficult at first
because all functions are anonymous---there is no name that a function can use to refer to itself.
Surprisingly, however, it is possible to make functions that behave recursively. Roughly, our
strategy will build on the ``self-replicating'' behavior of $\Omega$ to let functions invoke copies
of themselves an unbounded number of times.

Let's consider how we would like to write a factorial function:
\[
  \FACT\triangleq\lambda n.\,\IF\ (\ISZERO\ n)\ 1\ (\TIMES\ n\ (\FACT\ (\PRED\ n)))
\]
In slightly more readable notation, this is just:
\[
  \FACT\triangleq\lambda n.\ \ite{n=0}{1}{n\times\FACT\ (n-1)}
\]
Here, as in the definition above, the name $\FACT$ is simply meant to be shorthand for the
expression on the right-hand side of the equation. But $\FACT$ appears on the right-hand side of the
equation as well! Trying to expand the definition would yield an infinite term. This is not a
definition, it's a recursive equation.

\subsection{Recursion Removal Trick}

We can perform a ``trick'' to define a function $\FACT$ that satisfies the recursive equation above.
First, let's define a new function $\FACT'$ that looks like $\FACT$, but takes an additional
argument $f$. We assume that the function $f$ will be instantiated with $\FACT'$ itself.
\[
  \FACT'\triangleq\lambda f.\,\lambda n.\ \ite{n=0}{1}{n\times(f\ f\ (n-1))}
\]
Note that when we call $f$, we pass it a copy of itself, preserving the assumption that the actual
argument for $f$ will be $\FACT'$. Now we can define the factorial function $\FACT$ in terms of
$\FACT'$.
\[
  \FACT\triangleq\FACT'\ \FACT'
\]
Let's try evaluating $\FACT$ on an integer.
\begin{align*}
  \FACT\ 3&=(\FACT'\ \FACT')\ 3 && \text{Definition of } \FACT\\
  &=((\lambda f.\,\lambda n.\ \ite{n=0}{1}{n\times(f\ f\ (n-1))})\ \FACT')\ 3 && \text{Definition of } \FACT'\\
  &\to(\lambda n.\ \ite{n=0}{1}{n\times(\FACT'\ \FACT'\ (n-1))})\ 3 && \text{Application to } \FACT'\\
  &\to\ite{3=0}{1}{3\times(\FACT'\ \FACT'\ (3-1))} && \text{Application to } n\\
  &\to 3\times(\FACT'\ \FACT'\ (3-1)) && \text{Evaluating if}\\
  &\to^*\ldots\\
  &\to^* 3\times 2\times 1\times 1\\
  &\to^* 6
\end{align*}
So we now have a technique for writing a recursive function $f$: write a function $f'$ that
explicitly takes a copy of itself as an argument, and then define $f\triangleq f'\ f'$.

\subsection{Fixed point combinators}

There is another way of writing recursive functions: we can express the recursive function as the
fixed point of some other, higher-order function, and then take its fixed point. We saw this
technique earlier in the course when we defined the denotational semantics for while loops.

Let's consider the factorial function again. The factorial function $\FACT$ is a fixed point of the
following function.
\[
  G\triangleq\lambda f.\,\lambda n.\ \ite{n=0}{1}{n\times(f\ (n-1))}
\]
Recall that if $g$ is a fixed point of $G$, then we have $G\ g=g$. So if we had some way of finding
a fixed point of $G$, we would have a way of defining the factorial function $\FACT$.

There are a number of ``fixed point combinators,'' and the (infamous) $\Y$ combinator is one of
them. Thus, we can define the factorial function $\FACT$ to be simply $\Y\ G$, the fixed point of
$G$. The $\Y$ combinator is defined as
\[
  \Y\triangleq\lambda f.\,(\lambda x.\,f\ (x\ x))\ (\lambda x.\,f\ (x\ x)).
\]
It was discovered by Haskell Curry, and is one of the simplest fixed-point combinators. Note how
similar its definition is to $\Omega$.

Unfortunately, $\Y$ does not work under CBV. Evaluating $\Y\ G$ gives
\begin{align*}
  \Y\ G&\to(\lambda x.\,G\ (x\ x))\ (\lambda x.\,G\ (x\ x))\\
  &\to G\ ((\lambda x.\,G\ (x\ x))\ (\lambda x.\,G\ (x\ x)))\\
  &\to G\ (G\ ((\lambda x.\,G\ (x\ x))\ (\lambda x.\,G\ (x\ x))))\\
  &\to\cdots
\end{align*}
The argument to $G$ is never a value, so evaluation never terminates. Instead, we'll use a slight
variant of the $\Y$ combinator, called $\Z$, which works under CBV. It wraps the self-application
$x\ x$ in a $\lambda$, which delays it until it is needed. The $\Z$ combinator is defined as
\[
  \Z\triangleq\lambda f.\,(\lambda x.\,f\ (\lambda y.\,x\ x\ y))\ (\lambda x.\,f\ (\lambda y.\,x\ x\ y))
\]
Let's see it in action, on our function $G$. Define $\FACT$ to be $\Z\ G$ and calculate as follows:
\begin{align*}
  &\FACT\\
  &=\Z\ G\\
  &=(\lambda f.\,(\lambda x.\,f\ (\lambda y.\,x\ x\ y))\ (\lambda x.\,f\ (\lambda y.\,x\ x\ y)))\ G
    && \text{Definition of } \Z\\
  &\to(\lambda x.\,G\ (\lambda y.\,x\ x\ y))\ (\lambda x.\,G\ (\lambda y.\,x\ x\ y))\\
  &\to G\ (\lambda y.\,(\lambda x.\,G\ (\lambda y.\,x\ x\ y))\ (\lambda x.\,G\ (\lambda y.\,x\ x\ y))\ y)\\
  &=(\lambda f.\,\lambda n.\ \ite{n=0}{1}{n\times(f\ (n-1))})\\
  &\qquad(\lambda y.\,(\lambda x.\,G\ (\lambda y.\,x\ x\ y))\ (\lambda x.\,G\ (\lambda y.\,x\ x\ y))\ y)
    && \text{definition of } G\\
  &\to\lambda n.\ \ite{n=0}{1}{}\\
  &\qquad\mathbf{else}\ n\times((\lambda y.\,(\lambda x.\,G\ (\lambda y.\,x\ x\ y))\ (\lambda x.\,G\ (\lambda y.\,x\ x\ y))\ y)\ (n-1))\\
  &=_\beta\lambda n.\ \ite{n=0}{1}{n\times(\lambda y.\,(\Z\ G)\ y)\ (n-1)}\\
  &=_\beta\lambda n.\ \ite{n=0}{1}{n\times(\Z\ G\ (n-1))}\\
  &=\lambda n.\ \ite{n=0}{1}{n\times(\FACT\ (n-1))}
\end{align*}
There are many (indeed infinitely many) fixed-point combinators. Here's a cute one, due to Jan Willem
Klop:
\[
  \Y_k\triangleq(L\ L\ L\ L\ L\ L\ L\ L\ L\ L\ L\ L\ L\ L\ L\ L\ L\ L\ L\ L\ L\ L\ L\ L\ L\ L)
\]
where
\[
  L\triangleq\lambda a\,b\,c\,d\,e\,f\,g\,h\,i\,j\,k\,l\,m\,n\,o\,p\,q\,s\,t\,u\,v\,w\,x\,y\,z\,r.\
  (r\ (t\ h\ i\ s\ i\ s\ a\ f\ i\ x\ e\ d\ p\ o\ i\ n\ t\ c\ o\ m\ b\ i\ n\ a\ t\ o\ r))
\]
To gain some more intuition for fixed-point combinators, let's derive a fixed-point combinator that
was originally discovered by Alan Turing~\cite{turing}. Suppose we have a higher order function $f$,
and want the fixed point of $f$. We want $\Theta\ f$ to be a fixed point of $f$, so we need
\[
  \Theta\ f=f\ (\Theta\ f).
\]
This means, that we can write the following recursive equation:
\[
  \Theta=\lambda f.\,f\ (\Theta\ f).
\]
Now we can use the recursion removal trick we described earlier. Define $\Theta'$ as
$\lambda t.\,\lambda f.\,f\ (t\ t\ f)$, and $\Theta$ as $\Theta'\ \Theta'$. Then we have the
following equalities:
\begin{align*}
  \Theta&=\Theta'\ \Theta'\\
  &=(\lambda t.\,\lambda f.\,f\ (t\ t\ f))\ \Theta'\\
  &\to\lambda f.\,f\ (\Theta'\ \Theta'\ f)\\
  &=\lambda f.\,f\ (\Theta\ f)
\end{align*}
Let's try out the Turing combinator on our higher order function $G$ that we used to define
$\FACT$. This time we will use CBN evaluation.
\begin{align*}
  \FACT&=\Theta\ G\\
  &=((\lambda t.\,\lambda f.\,f\ (t\ t\ f))\ (\lambda t.\,\lambda f.\,f\ (t\ t\ f)))\ G\\
  &\to(\lambda f.\,f\ ((\lambda t.\,\lambda f.\,f\ (t\ t\ f))\ (\lambda t.\,\lambda f.\,f\ (t\ t\ f))\ f))\ G\\
  &\to G\ ((\lambda t.\,\lambda f.\,f\ (t\ t\ f))\ (\lambda t.\,\lambda f.\,f\ (t\ t\ f))\ G)\\
  &=G\ (\Theta\ G) && \text{for brevity}\\
  &=(\lambda f.\,\lambda n.\ \ite{n=0}{1}{n\times(f\ (n-1))})\ (\Theta\ G) && \text{Definition of } G\\
  &\to\lambda n.\ \ite{n=0}{1}{n\times((\Theta\ G)\ (n-1))}\\
  &=\lambda n.\ \ite{n=0}{1}{n\times(\FACT\ (n-1))}
\end{align*}
Note that the step from $G\ (\Theta\ G)$ passes the non-value $\Theta\ G$ to $G$, which is only
allowed under CBN. Under CBV, $\Theta\ G$ diverges, just like $\Y\ G$. As with $\Z$, we can fix this
by delaying the self-application, using $\Theta'_v\triangleq\lambda t.\,\lambda f.\,f\ (\lambda y.\,t\ t\ f\ y)$
in place of $\Theta'$.
% \end{document}
